\chapter{Unit 1: Introduction to Environment, Sustainability Metrics \& SDGs}

\begin{formulabox}{Atmospheric, Lithospheric \& Hydrospheric Invariants}
\begin{itemize}[leftmargin=1.2em]
    \item \textbf{Four Environmental Spheres}:
    \begin{itemize}
        \item \textbf{Atmosphere}: Total mass $\approx 5.15 \times 10^{18}\text{ kg}$. Dry air composition by volume: $N_2 = 78.08\%$, $O_2 = 20.95\%$, $Ar = 0.93\%$, $CO_2 = 0.042\%$ ($420\text{ ppm}$), trace gases.
        \item \textbf{Atmospheric Layers}:
        \begin{itemize}
            \item \textbf{Troposphere} ($0\text{--}12\text{ km}$): Contains $80\%$ atmospheric mass; temperature decreases at Environmental Lapse Rate (ELR $\approx 6.5^\circ\text{C/km}$).
            \item \textbf{Stratosphere} ($12\text{--}50\text{ km}$): Temperature inversion due to ozone UV absorption; Peak ozone layer at $20\text{--}30\text{ km}$.
            \item \textbf{Mesosphere} ($50\text{--}85\text{ km}$): Coldest layer (down to $-90^\circ\text{C}$); meteor ablation.
            \item \textbf{Thermosphere} ($85\text{--}600\text{ km}$): High kinetic temperatures from solar X-ray/EUV photoionization.
        \end{itemize}
        \item \textbf{Hydrosphere}: Global water volume $\approx 1.386 \times 10^9\text{ km}^3$. Distribution: Saline ocean water $= 97.5\%$, Freshwater $= 2.5\%$.
        \item \textbf{Freshwater Partitioning}: Glaciers and polar ice caps $= 68.7\%$, Groundwater $= 30.1\%$, Surface water (lakes, rivers, swamps) $= 1.2\%$.
        \item \textbf{Lithosphere}: Continental crust (granitic, $30\text{--}50\text{ km}$ thick, density $\approx 2.7\text{ g/cm}^3$) vs Oceanic crust (basaltic, $5\text{--}10\text{ km}$ thick, density $\approx 3.0\text{ g/cm}^3$).
    \end{itemize}
\end{itemize}
\end{formulabox}

\begin{formulabox}{Environmental Impact \& Carrying Capacity Formulas}
\begin{itemize}[leftmargin=1.2em]
    \item \textbf{Ehrlich-Holdren IPAT Identity}:
    \begin{equation}
    I = P \times A \times T
    \end{equation}
    where $I = \text{Environmental Impact}$, $P = \text{Population size}$, $A = \text{Affluence (GDP or consumption per capita)}$, $T = \text{Technology (impact per unit consumption)}$.
    \item \textbf{ImPACT Identity}:
    \begin{equation}
    I = P \times A \times C \times T'
    \end{equation}
    where $C = \text{energy/resource consumption per unit GDP}$, and $T' = \text{pollution per unit energy consumed}$.
    \item \textbf{Ecological Footprint ($EF$) vs Biocapacity ($BC$)}:
    \begin{equation}
    \text{Ecological Deficit/Reserve} = BC - EF
    \end{equation}
    If $EF > BC$, humanity is in global ecological overshoot ($1.75$ Earth equivalents currently).
    \item \textbf{Brundtland Definition of Sustainable Development (1987)}:
    \textit{``Development that meets the needs of the present without compromising the ability of future generations to meet their own needs.''} (Report: \textit{Our Common Future}).
    \item \textbf{Three Pillars of Sustainability (Triple Bottom Line)}: Environmental (Planet), Social (People), and Economic (Profit).
    \item \textbf{Weak vs Strong Sustainability}:
    \begin{itemize}
        \item \textbf{Weak Sustainability}: Assumes natural capital ($K_n$) and manufactured human capital ($K_m$) are fully fungible substitutes ($K_{\text{total}} = K_n + K_m = \text{constant}$).
        \item \textbf{Strong Sustainability}: Demands non-substitutability; critical natural capital ($K_n^*$) cannot be replaced by manufactured physical or human capital.
    \end{itemize}
\end{itemize}
\end{formulabox}

\begin{lawbox}{The 17 UN Sustainable Development Goals (SDGs 2016--2030)}
\begin{itemize}[leftmargin=1.2em]
    \item \textbf{SDG 1}: No Poverty \quad \textbullet\ \textbf{SDG 2}: Zero Hunger \quad \textbullet\ \textbf{SDG 3}: Good Health and Well-being
    \item \textbf{SDG 4}: Quality Education \quad \textbullet\ \textbf{SDG 5}: Gender Equality
    \item \textbf{SDG 6}: \textbf{Clean Water and Sanitation} (Safe drinking water, wastewater recycling, IWRM).
    \item \textbf{SDG 7}: \textbf{Affordable and Clean Energy} (Doubling renewable energy share, energy efficiency).
    \item \textbf{SDG 8}: Decent Work and Economic Growth \quad \textbullet\ \textbf{SDG 9}: Industry, Innovation and Infrastructure
    \item \textbf{SDG 10}: Reduced Inequalities
    \item \textbf{SDG 11}: \textbf{Sustainable Cities and Communities} (Air quality, transit, municipal waste).
    \item \textbf{SDG 12}: \textbf{Responsible Consumption and Production} (Circular economy, 3R, food loss).
    \item \textbf{SDG 13}: \textbf{Climate Action} (Paris Agreement Nationally Determined Contributions - NDCs).
    \item \textbf{SDG 14}: \textbf{Life Below Water} (Marine conservation, halting ocean acidification and plastic pollution).
    \item \textbf{SDG 15}: \textbf{Life on Land} (Halt deforestation, combat desertification, halt biodiversity loss).
    \item \textbf{SDG 16}: Peace, Justice and Strong Institutions \quad \textbullet\ \textbf{SDG 17}: Partnerships for the Goals.
\end{itemize}
\end{lawbox}

\begin{memorytip}
Remember the key environmental SDGs with the mnemonic \textbf{``W-E-C-C-C-L-L''}: Water (6), Energy (7), Cities (11), Consumption (12), Climate (13), Life below water (14), Life on land (15)!
\end{memorytip}
